\documentclass[11pt,leqno]{article}
\headheight=13.6pt

% packages
\usepackage[alphabetic]{amsrefs}
\usepackage{physics}
% margin spacing
\usepackage[top=1in, bottom=1in, left=0.5in, right=0.5in]{geometry}
\usepackage{amsfonts, amsmath, amssymb, amsthm}
\usepackage{extarrows}
\usepackage[none]{hyphenat}
\usepackage{fancyhdr}
\usepackage{graphicx}
\graphicspath{{./images/}}
\usepackage{float}
\usepackage{color}
\newcommand{\sai}[1]{\textcolor{red}{#1}}
\usepackage{enumitem}
% \usepackage{mathrsfs}
\usepackage{hyperref}
\usepackage[noabbrev, capitalise]{cleveref}
\crefformat{equation}{equation~#2#1#3}
\crefformat{lemma}{\textrm{Lemma}~#2#1#3}
\usepackage{quiver}

% theorems
\theoremstyle{plain}
\newtheorem{lem}{Lemma}
\newtheorem{lemma}[lem]{Lemma}
\newtheorem{thm}[lem]{Theorem}
\newtheorem{theorem}[lem]{Theorem}
\newtheorem{prop}[lem]{Proposition}
\newtheorem{proposition}[lem]{Proposition}
\newtheorem{cor}[lem]{Corollary}
\newtheorem{corollary}[lem]{Corollary}
\newtheorem{conj}[lem]{Conjecture}
\newtheorem{fact}[lem]{Fact}
\newtheorem{form}[lem]{Formula}

\theoremstyle{definition}
\newtheorem{defn}[lem]{Definition}
\newtheorem{definition/}[lem]{Definition}
\newenvironment{definition}
  {\renewcommand{\qedsymbol}{\textdagger}%
   \pushQED{\qed}\begin{definition/}}
  {\popQED\end{definition/}}
\newtheorem{example}[lem]{Example}
\newtheorem{remark}[lem]{Remark}
\newtheorem{exercise}[lem]{Exercise}
\newtheorem{notation}[lem]{Notation}

\numberwithin{equation}{section}
\numberwithin{lem}{section}

% header/footer formatting
\pagestyle{fancy}
\fancyhead{}
\fancyfoot{}
\fancyhead[L]{Bezrukavnikov-Mirkovic-Rumynin localization for $\SL_2$}
\fancyhead[C]{}
\fancyhead[R]{Sai Sivakumar}
\fancyfoot[R]{\thepage}
\renewcommand{\headrulewidth}{1pt}

% paragraph indentation/spacing
\setlength{\parindent}{0cm}
\setlength{\parskip}{10pt}
\renewcommand{\baselinestretch}{1.25}

% mathscript S
\usepackage[mathscr]{euscript}

% math operators
\DeclareMathOperator{\Stab}{Stab}
\DeclareMathOperator{\Ind}{Ind}
\let\sl\undefined
\DeclareMathOperator{\sl}{\mathfrak{sl}}
\DeclareMathOperator{\SL}{SL}

% bracket notation for inner product
\usepackage{mathtools}

\DeclarePairedDelimiterX{\abr}[1]{\langle}{\rangle}{#1}

% smileys frownies
\usepackage{wasysym}
\newcommand{\smallhappy}{\raisebox{-.14em}{\smiley}}
\newcommand{\happy}{\raisebox{-.24em}{\resizebox{1.2em}{!}{\smiley}}}
\newcommand{\smallsad}{\raisebox{-.14em}{\frownie}}
\newcommand{\sad}{\raisebox{-.24em}{\resizebox{1.2em}{!}{\frownie}}}
\DeclareMathOperator{\mathhappy}{\!\happy\!}
\DeclareMathOperator{\smallmathhappy}{\!\smallhappy\!}
\DeclareMathOperator{\mathsad}{\!\sad\!}
\DeclareMathOperator{\smallmathsad}{\!\smallsad\!}

% set page count index to begin from 1
\setcounter{page}{1}

% array column and row separation
\arraycolsep = 2pt
\renewcommand{\arraystretch}{.6}

\begin{document}

These are calculations from from the 2026 FRC: Representation Theory at Louisiana State University. I was part of Working Group 2: Localization and representations of Lie algebras in positive characteristic, supervised by Pramod Achar and Pablo Boixeda Alvarez. Our TA was John O'Brien. The other students in this group were Charlie Barth, Merrick Cai, Qianyi Cao, Wade Cheng, Frank Wang, and Mike Young. Attached below is the abstract the organizers provided.
\hrulefill

The representation theory of complex semisimple Lie algebras has been studied since the beginning of the 20th century. In 1981, Beilinson-Bernstein and Brylinski-Kashiwara proved a breakthrough result, called the Localization Theorem, relating representations of complex semisimple Lie algebras to differential operators on flag manifolds and thus providing geometric tools to the study of these representations. The representation theory of Lie algebras in characteristic $p$ behaves very differently, but in some ways it is simpler. In particular, all irreducible representations are finite-dimensional. However, the naive analogue of the Localization Theorem is false. Nevertheless, in 2008, Bezrukavnikov-Mirkovic-Rumynin discovered and proved a modified version of the Localization Theorem that is valid in positive characteristic and thus provides similar geometric tools for understanding the representation theory in characteristic $p$. This course will be about this BMR localization theorem. A rough outline of topics to be covered is:
\begin{enumerate}[label = \roman*.]
  \item Introduction to Lie algebra representations in characteristic $p$
  \item Geometry of flag varieties
  \item Differential operators and $D$-modules in characteristic $p$
  \item Derived localization
  \item The Azumaya property and coherent sheaves
\end{enumerate}
We'll also work out hands-on examples for the Lie algebra $\sl_2$.

\hrulefill

The listed prerequisites were algebraic geometry (as in Hartshorne chapters 1-2) and representation theory of Lie algebras (as in Fulton and Harris).

Our main source is \href{https://annals.math.princeton.edu/wp-content/uploads/annals-v167-n3-p05.pdf}{Localization of modules for a semisimple
Lie algebra in prime characteristic} by Bezrukavnikov, Mirkovic, and Rumynin, but we also follow \sai{sources and links}

A \hyperref[glossary]{glossary} is attached at the end of the document.\newpage

\section*{Calculations for $\SL_2$}
Throughout, let $k$ be an algebraically closed field of characteristic $p>0$, $G = \SL_2$, fix $B$ the Borel subgroup of upper-triangular matrices, $T$ the subgroup of diagonal matrices, $N = [B,B]$ the subgroup of strictly upper-triangular matrices. Denote by $\mathfrak g$, $\mathfrak b$, $\mathfrak t$, $\mathfrak n$ the corresponding Lie algebras. Our preferred presentation of $\mathfrak g= \sl_2$ is $k\{e = \bigl(\!\begin{smallmatrix}
	0 & 1 \\ 0 & 0
\end{smallmatrix}\!\bigr), h = \bigl(\!\begin{smallmatrix}
	1 & 0 \\ 0 & -1
\end{smallmatrix}\!\bigr), f = \bigl(\!\begin{smallmatrix}
	0 & 0 \\ 1 & 0
\end{smallmatrix}\!\bigr)\}$, so that $\mathfrak b = k\{e,h\}$, $\mathfrak t = k\{h\}$, and $\mathfrak n = k\{e\}$.

We will introduce more notation as needed throughout, and we will indicate when certain constructions require $p$ to be large enough (i.e., larger than the Coxeter number of $\sl_2$, which is $2$).

\sai{other notations, what theorems we want to manifest}
\subsection{Representation theory of $\sl_2$}
\subsection{Sheaf of differential operators on flag variety}
The flag variety $X\coloneqq G/B$ may be identified with $\mathbb P^1$ via $gB\mapsto g[1:0]$; denote $[1:0]$ by $\infty$ and $[0:1$ by $0$. The flag variety $X$ may also be identified with the collection $\mathcal B$ of Borel subalgebras of $\mathfrak g$ via $gB\mapsto g\mathfrak b g^{-1}$. We will move between both identifications as needed.

The sheaf of (crystalline) differential operators $\mathcal D_X$ on $X$ on the affine open $U \coloneqq \{[z:1]\mid z\in k\}\subseteq \mathbb P^1$ containing $0$ is given by $k\abr{z,\partial_z}$ with relation $[\partial_z,z]=1$. On the affine open $V\coloneqq \{[1:w]\mid w\in k\}\subseteq \mathbb P^1$ containing $\infty$ is given by $k\abr{w,\partial_w}$ with relation $[\partial_w,w]=1$. From now on, we suppress indicating relations for non-commutative algebras. Glue these differential operators on $U\cap V$ via $w = z^{-1}$ and $\partial_w = -z^2\partial_z$ (this is because $\partial_{z^{-1}}(z^n) = -nz^{n+1}$). It follows that the global sections of $\mathcal D_X$ (expressed in the coordinate on $U$) is given by 
\[k\abr{z^2\partial_z, z\partial_z, \partial_z}.\]
Recall that since $G$ acts on $X = G/B$, we obtain an injection $U(\mathfrak g)\hookrightarrow \Gamma(X,\mathcal D_X)$; in view of this, note that 
\begin{align*}
    z^2\partial_z = -\partial_w &= e = \bigl(\!\begin{smallmatrix}
	0 & 1 \\ 0 & 0
\end{smallmatrix}\!\bigr)\\
    z\partial_z = -w\partial_w &= h/2 = \bigl(\!\begin{smallmatrix}
	1/2 & 0 \\ 0 & -1/2
\end{smallmatrix}\!\bigr)\\
\partial_z = -w^2\partial_w &= -f = \bigl(\!\begin{smallmatrix}
	0 & 0 \\ -1 & 0
\end{smallmatrix}\!\bigr).
\end{align*}

For $p>2$, use the non-degenerate bilinear form $\abr{-,-}\colon \abr{x,y}\mapsto \Tr(xy)/2$ (the Killing form) on $\sl_2$ to obtain an identification of $(\mathfrak g/\mathfrak b)^\ast$ with $\mathfrak b^\perp$. This gives rise to an isomorphism of the cotangent bundle $T^\ast X\to X$ with the associated bundle $G\times^B\mathfrak b^\perp\to \mathcal B$, via the map $(gB, \xi)\mapsto (g, \overline{\xi})$

\newpage
\section*{Glossary}
\label{glossary}
\end{document}